% Anforderungen und Bedrohungsmodell — claude-tunnel bachelor thesis.
\chapter{Anforderungen und Bedrohungsmodell}
\label{ch:anforderungen}

Dieses Kapitel praezisiert, was das System leisten soll
(Abschnitt~\ref{sec:anf-funktional} und~\ref{sec:anf-nichtfunktional}), gegen
wen es schuetzen soll (Abschnitt~\ref{sec:bedrohung}) und welche Schutzziele
daraus folgen (Abschnitt~\ref{sec:schutzziele}). Es bildet die Grundlage, an der
die Architektur (Kapitel~\ref{ch:anforderungen}) und die Evaluation gemessen
werden.

\section{Funktionale Anforderungen}
\label{sec:anf-funktional}

\begin{description}
  \item[F1 Tunnel beliebiger Protokolle.] Das System tunnelt bidirektionale
    \gls{ac:tcp}-Stroeme sowie \gls{ac:udp}-Datagramme zwischen einem Client und
    einem Origin. Datagrammgrenzen bleiben erhalten.
  \item[F2 Enrollment und Identitaet.] Ein Agent registriert sich ueber ein
    einmaliges Join-Token bei der Steuerebene und bindet dabei einen
    ed25519-Schluessel; kurzlebige Anmeldedaten dienen als mTLS-aehnlicher
    Nachweis gegenueber dem Edge.
  \item[F3 Rendezvous.] Ein Client loest ein opakes Routing-Token beim Edge zum
    zustaendigen Tunnel auf, ohne einen Hostnamen zu nennen.
  \item[F4 Missbrauchsschranke.] Die Rendezvous-Anbahnung ist durch einen
    \gls{ac:pow}-Nachweis verteuert; die Schwierigkeit ist konfigurierbar.
  \item[F5 Direkter Pfad und Fallbacks.] Bei Erreichbarkeit nutzt der Client
    einen direkten \gls{ac:p2p}-Pfad zum Agent; andernfalls den Edge-Relay. Ist
    \gls{ac:udp} blockiert, traegt ein \gls{ac:tls}-ueber-\gls{ac:tcp}-Fallback
    denselben Tunnel.
  \item[F6 Steuerebene.] Enrollment, Registry/Rendezvous und Abrechnung sind
    ueber einen laufenden Dienst erreichbar.
  \item[F7 Beobachtbarkeit.] Der Agent exponiert Prometheus-Metriken
    (Tunnelzahl, Bytes, Handshake-Latenz).
  \item[F8 Konten und Zahlung.] Ein pseudonymes Prepaid-Guthaben gated die
    Token-Ausgabe; das Guthaben wird ueber einen Krypto-Payment-Stub aufgeladen.
\end{description}

\section{Nicht-funktionale Anforderungen}
\label{sec:anf-nichtfunktional}

\begin{description}
  \item[N1 Providerblindheit.] Der Betreiber der Vermittlungsebene darf Ziel und
    Nutzlast strukturell nicht einsehen koennen; er relayt nur Chiffretext.
  \item[N2 Kein Vorhalten von Vertrauensmaterial.] Die Steuerebene haelt weder
    Origin-Schluessel noch Nutzlast (thin control plane).
  \item[N3 Beschraenkter Overhead.] Der Latenz-Overhead gegenueber der reinen
    Netzlaufzeit soll klein und -- ueber Netzbedingungen -- vorhersagbar sein
    (empirisch belegt in der Evaluation).
  \item[N4 Reproduzierbarkeit.] Aufbau, Betrieb und Messung erfolgen
    ausschliesslich containerbasiert und sind aus dem Repository reproduzierbar.
  \item[N5 Pseudonymitaet.] Konten werden ausschliesslich ueber opake Kennungen
    gefuehrt; es werden keine personenbezogenen Daten gespeichert.
\end{description}

\section{Bedrohungsmodell}
\label{sec:bedrohung}

\subsection{Akteure und Faehigkeiten}

\begin{description}
  \item[A1 Neugieriger Betreiber (honest-but-curious).] Der Edge/Betreiber
    fuehrt das Protokoll korrekt aus, versucht aber, aus dem vermittelten
    Verkehr Ziel oder Inhalt zu lernen. Er sieht Quell-/Zieladressen der
    QUIC-Pakete, Zeitpunkte und Volumina, jedoch -- so das Ziel -- niemals
    Klartext oder Ziel-Hostnamen.
  \item[A2 Aktiver Netz-Angreifer / Zensor.] Kann Pakete beobachten, verwerfen
    und Protokolle grob blockieren (insbesondere \gls{ac:udp}). Kann versuchen,
    sich als Origin auszugeben.
  \item[A3 Boesartiger Client / Fluter.] Versucht, ueber massenhafte
    Rendezvous-Versuche Ressourcen zu erschoepfen.
  \item[A4 Finanzierter Sybil-Angreifer.] Kann beliebig viele pseudonyme Konten
    eroeffnen und bezahlen.
\end{description}

\subsection{Vertrauensgrenzen und Annahmen}

Der Client vertraut der \emph{Origin-Identitaet}, die er ausserhalb des Bandes
als Teil einer Capability erhaelt (Pinning, vgl.
Abschnitt~\ref{sec:noise}). Der Agent ist \emph{Verwalter} des
Origin-Noise-Schluessels; dieser verlaesst den Agent nie. Die Steuerebene wird
als verfuegbar, aber nicht als vertraulich angenommen (N2). Die kryptographischen
Primitiven (X25519, ChaCha20-Poly1305, BLAKE2s, ed25519, SHA-256) gelten als
sicher.

\section{Schutzziele und Nicht-Ziele}
\label{sec:schutzziele}

\paragraph{Schutzziele.}
\textbf{S1 Ende-zu-Ende-Vertraulichkeit und -Integritaet} der Nutzlast gegenueber
A1 und A2 (durch Noise-Ende-zu-Ende, Abschnitt~\ref{sec:noise}).
\textbf{S2 Ziel-Blindheit} des Betreibers gegenueber A1 (Adressierung ueber opake
Tokens, N1). \textbf{S3 Origin-Authentizitaet} gegenueber A2 (Pinning schlaegt
fehl bei falschem Schluessel). \textbf{S4 Missbrauchsdaempfung} gegenueber A3
(asymmetrischer \gls{ac:pow}-Aufwand).

\paragraph{Nicht-Ziele.}
Ausdruecklich \emph{nicht} adressiert werden: \emph{Anonymitaet} im Sinne der
Unverkettbarkeit von Client und Ziel gegenueber einem globalen Beobachter (das
ist Tors Ziel, vgl.\ Abschnitt~\ref{sec:rw-tor}); \emph{Traffic-Analyse- und
Verschleierungs-Resistenz} gegen A2 (Groessen- und Timing-Seitenkanaele bleiben
sichtbar); \emph{vollstaendiger DoS-Schutz} (\gls{ac:pow} daempft, verhindert
aber nicht); und die \emph{oekonomische Abwehr von A4}: ein finanzierter
Sybil-Angreifer kann beliebig viele Konten kaufen. Das prepaid-basierte Gating
verteuert Missbrauch, loest die Sybil-Oekonomie aber nicht; diese Grenze wird
offen benannt und im Backlog gefuehrt.
